% Part: incompleteness
% Chapter: representability-in-q
% Section: basic-representable

\documentclass[../../../include/open-logic-section]{subfiles}

\begin{document}

\olfileid{inc}{req}{bre}
\olsection{بنیادی فنکشن~$\Th{Q}$ وچ قابلِ نمائندگی نیں}

پہلاں سانوں وکھانا اے کہ سارے بنیادی فنکشن~$\Th{Q}$ وچ
قابلِ نمائندگی نیں۔ ایس لئی سانوں وکھانا اے کہ ہر $k$-دلیلی
بنیادی فنکشن $f(x_0,\dots,x_{k-1})$ نال !!a{formula}
$!A_f(x_0,\dots,x_{k-1},y)$ کِویں جوڑیے جیہڑا اوہدی نمائندگی کرے۔

اسیں صفر، جانشین، جمع، ضرب، برابری دے امتیازی فنکشن تے
پروجیکشن فنکشناں دی نمائندگی کر سکاں گے۔ ہر صورت وچ نمائندگی
کرن والا فارمولا سِدھا اے۔ مثلاً صفر دی نمائندگی فارمولا
$y = \Obj 0$ کردا اے، جانشین دی نمائندگی !!{formula} $x_0' = y$
کردا اے، تے جمع دی نمائندگی !!{formula} $(x_0 + x_1) = y$
کردا اے۔ اصل کم ایہ وکھانا اے کہ $\Th{Q}$ متعلقہ !!{sentence}s
ثابت کر سکدا اے۔ مثلاً اُتلا !!{formula} جمع دی نمائندگی
کردا اے جے ایہ وکھائیے کہ قدرتی عدداں $m$ تے $n$ دے ہر جوڑے
لئی $\Th{Q}$ ایہ دونوں گلاں ثابت کردا اے۔
\begin{align*}
& \eq[\num n + \num m][\num {n+m}] \text{ تے}\\
& \lforall[y][(\eq[(\num n + \num m)][y] \lif \eq[y][\num{n+m}])].
\end{align*}

\begin{prop}
\ollabel{prop:rep-zero}
صفر فنکشن $\Zero(x) = 0$ دی~$\Th{Q}$ وچ نمائندگی فارمولا
$!A_{\Zero}(x,y) \ident \eq[y][\Obj 0]$ کردا اے۔
\end{prop}

\begin{prop}
\ollabel{prop:rep-succ}
جانشین فنکشن $\Succ(x) = x+1$ دی~$\Th{Q}$ وچ نمائندگی فارمولا
$!A_{\Succ}(x,y) \ident \eq[y][x']$ کردا اے۔
\end{prop}

\begin{prop}
\ollabel{prop:rep-proj}
پروجیکشن فنکشن $\Proj{n}{i}(x_0, \dots, x_{n-1}) = x_i$
دی~$\Th{Q}$ وچ نمائندگی ہِٹھلا فارمولا کردا اے۔
\[!A_{\Proj{n}{i}}(x_0, \dots, x_{n-1}, y) \ident \eq[y][x_i].\]
\end{prop}

\begin{prob}
ثابت کرو کہ $\eq[y][\Obj 0]$، $\eq[y][x']$ تے $\eq[y][x_i]$ بالترتیب
$\Zero$، $\Succ$ تے $\Proj{n}{i}$ دی نمائندگی کردے نیں۔
\end{prob}

\begin{prop}
\ollabel{prop:rep-id}
% PNB-CMP-054: case دی عام عبارت لئی اک text wrapper؛فارمولے قائم۔
برابری~$=$ دا امتیازی فنکشن،
\[
\Char{=}(x_0, x_1) =
\begin{cases}
  1 & \text{جے } x_0 =x_1\\
  0 & \text{نہیں تاں}
\end{cases}
\]
دی~$\Th{Q}$ وچ نمائندگی ہِٹھلا فارمولا کردا اے۔
\[
  !A_{\Char{=}}(x_0, x_1, y) \ident (\eq[x_0][x_1] \land \eq[y][\num{1}]) \lor (\eq/[x_0][x_1] \land
\eq[y][\num{0}]).
\]
\end{prop}

ثبوت لئی ہِٹھلا معاون قضیہ چاہیدا اے۔

\begin{lem}
\ollabel{lem:q-proves-neq} قدرتی عدداں $n$ تے $m$ لئی جے $n \neq m$ ہووے، تاں $\Th{Q} \Proves \eq/[\num n][\num m]$۔
\end{lem}

\begin{proof}
$n$ اُتے استقرا نال وکھاؤ کہ ہر $m$ لئی، جے $n \neq m$، تاں $Q \Proves \eq/[\num n][\num m]$۔

بنیادی صورت وچ $n = 0$ اے۔ جے $m$، $0$ دے برابر نہیں، تاں
$m = k + 1$، جتھے $k$ کوئی قدرتی عدد اے۔ ساڈے کول اک مسلّمہ اے جیہڑا کہندا اے
$\lforall[x][\eq/[0][x']]$۔ مسوّر دے اک مسلّمے نال، $x$ دی تھاں $\num k$ رکھ کے، اسیں $\eq/[0][\num k']$ کڈھ سکدے آں۔
پر $\num k'$ بس $\num m$ ای اے۔

استقرائی قدم وچ اسیں $n$ لئی دعویٰ درست فرض کر سکدے آں،
تے $n+1$ نوں ویکھدے آں۔ $m$ کوئی وی قدرتی عدد ہووے۔ دو
صورتاں نیں:یا $m = 0$،یا کسے $k$ لئی $m = k+1$۔ پہلی صورت دا ثبوت اُتلے طریقے نال ہو جاندا اے۔ دوجی صورت وچ
فرض کرو $n+1 \neq
k+1$۔ تاں $n \neq k$۔ $n$ لئی استقرائی مفروضے نال ساڈے
کول $\Th{Q} \Proves \eq/[\num n][\num k]$ اے۔ ساڈے کول
اک مسلّمہ اے جیہڑا کہندا اے
$\lforall[x][\lforall[y][\eq[x'][y'] \lif \eq[x][y]]]$۔ مسوّر دا اک مسلّمہ ورت کے ساڈے کول
$\eq[\num n'][\num k'] \lif \eq[\num n][\num
  k]$ اے۔ قضیاتی منطق نال اسیں $\Th{Q}$ وچ ایہ کڈھ سکدے آں:
$\eq/[\num n][\num k] \lif \eq/[\num n'][\num k']$۔
اثباتِ مقدم دے قاعدے نال اسیں $\eq/[\num n'][\num k']$ کڈھ سکدے آں؛
ایہو سانوں چاہیدا اے، کیوں جے $\num k'$،$\num m$ اے۔
\end{proof}

\begin{explain}
دھیان دیو کہ معاون قضیہ بنیادی گل ایہ آکھدا اے کہ $\Th{Q}$ ثابت کر سکدا اے کہ وکھریاں
عددی عبارتاں وکھریاں شےواں دسدیاں نیں۔ مثال لئی،$\Th{Q}$
ایہ ثابت کردا اے: $0'' \neq 0'''$۔ پر ایہ دکھاؤن لئی کہ
گل ہر صورت وچ درست اے،ذرا احتیاط چاہیدی اے۔ ایہ وی دھیان رکھو کہ اسیں جیہڑا استقرا ورت رہے آں، اوہ
$\Th{Q}$ دے \emph{باہر} اے۔
\end{explain}

\begin{proof}[\olref{prop:rep-id} دا ثبوت]
جے $n = m$ ہووے، تاں $\num{n}$ تے $\num{m}$ اکّو ٹرم نیں،
تے $\Char{=}(n, m) = 1$۔ پر
$\Th{Q} \Proves (\eq[\num{n}][\num{m}] \land
\eq[\num{1}][\num{1}])$،سو اوہ
$!A_=(\num{n}, \num{m},
\num{1})$ ثابت کردا اے۔ جے $n \neq m$ ہووے، تاں
$\Char=(n, m) = 0$۔ \olref{lem:q-proves-neq} نال،
$\Th{Q} \Proves \eq/[\num{n}][\num{m}]$،تے ایس لئی
$(\eq/[\num{n}][\num{m}] \land \Obj 0 = \Obj 0)$ وی۔
سو $\Th{Q}
\Proves !A_=(\num{n}, \num{m}, \num{0})$۔

دوجے حصے لئی وی دو صورتاں نیں۔ جے $n = m$،سانوں
دکھانا اے کہ $\Th{Q} \Proves \lforall[y][(!A_=(\num{n}, \num{m}, y)
  \lif \eq[y][\num{1}])]$۔ غیر رسمی دلیل وچ فرض کرو
$!A_=(\num{n},
\num{m}, y)$،یعنی:
\[
(\eq[\num{n}][\num{n}] \land \eq[y][\num{1}]) \lor
(\eq/[\num{n}][\num{n}] \land \eq[y][\num{0}])
\]
فصل دا کھبا جز منطق نال $\eq[y][\num{1}]$ دا نتیجہ دیندا اے؛
اوہدا سجا جز $\eq[\num{n}][\num{n}]$ نال ٹکراندا اے،جیہڑا
منطق نال ثابت ہو سکدا اے۔

دوجی صورت وچ فرض کرو $n \neq m$۔ تاں
$!A_=(\num{n},
\num{m}, y)$ ایہ اے:
\[
(\eq[\num{n}][\num{m}] \land \eq[y][\num{1}]) \lor
(\eq/[\num{n}][\num{m}] \land \eq[y][\num{0}])
\]
ایتھے فصل دا کھبا جز $\eq/[\num{n}][\num{m}]$ نال ٹکراندا اے،
جیہڑا \olref{lem:q-proves-neq} نال $\Th{Q}$ وچ ثابت ہو سکدا اے؛
فصل دا سجا جز $\eq[y][\num{0}]$ دا نتیجہ دیندا اے۔
\end{proof}

\begin{prop}
\ollabel{prop:rep-add}
جمع فنکشن $\Add(x_0, x_1) = x_0+x_1$ دی~$\Th{Q}$ وچ نمائندگی ہِٹھلا فارمولا کردا اے۔
\[
  !A_{\Add}(x_0, x_1, y) \ident \eq[y][(x_0 + x_1)].
\]
\end{prop}

\begin{lem}
\ollabel{lem:q-proves-add}
$\Th{Q} \Proves \eq[(\num{n} + \num{m})][\num{n+m}]$
\end{lem}

\begin{proof}
اسیں ایہ $m$ اُتے استقرا نال ثابت کردے آں۔ جے $m = 0$،
دعویٰ ایہ اے کہ
$\Th{Q} \Proves \eq[(\num{n} + \Obj 0)][\num{n}]$۔ ایہ
مسلّمے~$!Q_4$ توں نکلدا اے۔ ہن $m$ لئی دعویٰ فرض کرو؛
ہن $m+1$ لئی دعویٰ ثابت کریے۔ یعنی وکھائیے کہ
$\Th{Q} \Proves \eq[(\num{n} +
  \num{m+1})][\num{n+m+1}]$۔ دھیان دیو کہ $\num{m+1}$
بس $\num{m}'$ اے، تے $\num{n+m+1}$ بس $\num{n+m}'$ اے۔
مسلّمے $!Q_5$ نال،$\Th{Q}
\Proves \eq[(\num{n} + \num{m}')][(\num{n}+\num{m})']$۔
استقرائی مفروضے نال،
$\Th{Q} \Proves \eq[(\num{n} + \num{m})][\num{n+m}]$۔ سو
$\Th{Q} \Proves \eq[(\num{n} + \num{m}')][\num{n+m}']$۔
\end{proof}

\begin{proof}[\olref{prop:rep-add} دا ثبوت]
!!{formula}~$!A_\Add(x_0, x_1, y)$ جیہڑا $\Add$ دی نمائندگی کردا اے، ایہ اے
$\eq[y][(x_0 + x_1)]$۔ پہلاں اسیں دکھاندے آں کہ جے
$\Add(n, m) = k$،تاں
$\Th{Q} \Proves !A_\Add(\num{n}, \num{m}, \num{k})$،
یعنی $\Th{Q}
\Proves \eq[\num{k}][(\num{n} + \num{m})]$۔ پر کیوں جے
$k = n + m$، ایس لئی $\num{k}$ عین $\num{n+m}$ اے،تے اسیں
\olref{lem:q-proves-add} وچ ایہ دکھا چُکے آں:
$\Th{Q} \Proves \eq[(\num{n} +
  \num{m})][\num{n+m}]$۔

سانوں ایہ وی دکھانا اے کہ جے $\Add(n, m) = k$،تاں
\[
\Th{Q} \Proves \lforall[y][(!A_\Add(\num{n}, \num{m}, y) \lif
  \eq[y][\num{k}])].
\]
فرض کرو ساڈے کول $\eq[(\num{n} + \num{m})][y]$ اے۔
کیوں جے
\[
\Th{Q} \Proves \eq[(\num{n}+\num{m})][\num{n+m}],
\]
اسیں کھبے پاسے دی تھاں $\num{n+m}$ رکھ کے
$\eq[\num{n+m}][y]$ حاصل کر سکدے آں،ہر~$y$ لئی۔
\end{proof}

\begin{prop}
\ollabel{prop:rep-mult}
ضرب فنکشن $\Mult(x_0, x_1) = x_0 \cdot x_1$ دی~$\Th{Q}$
وچ نمائندگی ہِٹھلا فارمولا کردا اے۔
\[
  !A_{\Mult}(x_0, x_1, y) \ident y = (x_0 \times x_1).
\]
\end{prop}

\begin{proof} مشق۔ \end{proof}

\begin{lem}
\ollabel{lem:q-proves-mult}
$\Th{Q} \Proves \eq[(\num{n} \times \num{m})][\num{n \cdot m}]$
\end{lem}

\begin{proof} مشق۔ \end{proof}

\begin{prob}
\olref[inc][req][bre]{lem:q-proves-mult} ثابت کرو۔
\end{prob}

\begin{prob}
\olref[inc][req][bre]{lem:q-proves-mult} ورت کے
\olref[inc][req][bre]{prop:rep-mult} ثابت کرو۔
\end{prob}

\begin{explain}
  یاد رکھو کہ اسیں حساب دی زبان دے فنکشن نشان لئی $\times$
ورتدے آں تے عدداں اُتے عام ضرب دے عمل لئی $\cdot$۔ ایس لئی
$\cdot$ عدد دسن والیاں عبارتاں وچکار آ سکدا اے، جیویں
$m \cdot n$ وچ۔ $\times$ صرف حساب دی زبان دے ٹرماں وچکار
آندا اے، جیویں $(\num{m} \times \num{n})$ وچ۔ ہور وی اُلجھن
والی گل ایہ اے کہ $+$، !!{function} تے جمع دے عمل، دونوں لئی
ورتیا جاندا اے۔ جدوں ایہ ٹرماں وچکار آوے، جیویں
$(\num{n} + \num{m})$ وچ، تاں ایہ حساب دی زبان دا $2$-دلیلی
!!{function} اے۔ جدوں ایہ عدداں وچکار آوے، جیویں $n+m$ وچ،
تاں ایہ جمع دا عمل اے۔ $\num{n+m}$ وچ وی ایہو جمع اے۔
ایہ عدد~$n+m$ دی معیاری عددی عبارت اے۔
\end{explain}

\end{document}
